% !TeX root = ../本科微分几何笔记.tex
% 由主文档通过 \input 载入；请编译主文档。
\setcounter{chapter}{10} % 保留原书章号。
\chapter{从曲面到流形与黎曼几何}
\label{chap:dg-manifolds}

\section{抽象流形：保留坐标兼容性}
曲面的许多计算只用到二维局部坐标及坐标过渡，而不再需要空间位置 $\mathbf x$。这提示我们把局部欧氏性质作为出发点，允许研究不预先嵌入某个欧氏空间的几何对象。

\begin{definition}[光滑流形]
一个 $n$ 维\emph{光滑流形（smooth manifold）}是满足 Hausdorff 性与第二可数性的拓扑空间 $M$，配备一族覆盖它的局部坐标图 $\varphi:U\to\R^n$ 的开集，使每个图都是同胚，且交叠处的坐标过渡 $\psi\circ\varphi^{-1}$ 光滑。所有与这些图光滑相容的坐标图构成其光滑结构。
\end{definition}
Hausdorff 性保证不同点可用不相交邻域分离；第二可数性要求存在可数拓扑基，排除局部欧氏但整体过大的病态空间。曲面是 $n=2$ 的实例，圆周是 $n=1$ 的实例。有边界流形的图允许取值于半空间的相对开集，映到边界超平面的点构成流形边界；这不同于嵌入曲面区域在 $\R^3$ 中的普通集合边界。

\section{切向量、余切向量与张量}
抽象流形上没有可直接相减的位置向量。可把经过 $p$ 的曲线按“在某个局部坐标下具有同样速度”分成等价类，得到切空间 $T_pM$；链式法则保证该关系与坐标无关。在坐标 $u^1,\ldots,u^n$ 下，基向量写为 $\partial_i=\partial/\partial u^i$。

\emph{余切空间（cotangent space）} $T_p^*M$ 是 $T_pM$ 的对偶空间。基一形式 $\dd u^i$ 满足 $\dd u^i(\partial_j)=\delta^i_j$。向量分量与一形式分量变换方向相反；例如 $X(f)=\dd f(X)=\sum_iX^i\partial_i f$ 的数值与坐标无关。对切向量和余切向量多线性作用的对象称为\emph{张量（tensor）}；度量是一个对称的二阶协变张量，其坐标变换正是式~\eqref{eq:dg-metric-transform}。

\section{黎曼度量、梯度与 Laplace--Beltrami 算子}
\begin{definition}[黎曼度量]
若每个切空间 $T_pM$ 上给定一个内积 $g_p$，其坐标系数随 $p$ 光滑变化，则称 $g$ 为\emph{黎曼度量（Riemannian metric）}，$(M,g)$ 称为黎曼流形。长度、角度、面积或体积以及距离均由该度量定义。
\end{definition}
度量允许把一形式转成向量。对光滑函数 $f$，\emph{梯度（gradient）}定义为满足 $g(\grad f,X)=X(f)$ 的唯一向量场，因此
\[
 \grad f=\sum_{i,j}g^{ij}(\partial_jf)\partial_i.
\]
在坐标下直接把偏导数列表当作梯度分量，只在正交单位坐标中成立。度量还唯一确定满足无挠性与度量相容的\emph{Levi--Civita 联络（Levi--Civita connection）}；其系数仍由式~\eqref{eq:dg-christoffel} 给出，只把指标范围改为 $1,\ldots,n$。唯一性来自同样的度量求导组合，存在性可由该公式及其坐标变换性质验证。

记 $|g|=\det(g_{ij})$，体积元为 $\sqrt{|g|}\dd u^1\cdots\dd u^n$。向量场 $X$ 的散度坐标式为
\[
 \operatorname{div}X=\frac1{\sqrt{|g|}}\sum_i\partial_i(\sqrt{|g|}X^i).
\]
定义\emph{Laplace--Beltrami 算子（Laplace--Beltrami operator）}为 $\Delta f=\operatorname{div}(\grad f)$，从而
\begin{equation}
 \Delta f=\frac1{\sqrt{|g|}}\sum_{i,j}
 \partial_i\left(\sqrt{|g|}g^{ij}\partial_jf\right).
 \label{eq:dg-laplacian}
\end{equation}
本书取与欧氏 $\sum_i\partial_i^2$ 一致的符号，另一常见约定相差负号。二维保角度量 $e^{2\varphi}(\dd u^2+\dd v^2)$ 下，$\Delta f=e^{-2\varphi}(f_{uu}+f_{vv})$；梯度、散度和曲率都要使用度量，不能从坐标形式的相似性推断几何等同。

\section{三个抽象模型与后续方向}
欧氏平面度量 $\dd x^2+\dd y^2$ 的曲率为零；半径 $R$ 球面的曲率为 $1/R^2$；上半平面度量 $(\dd x^2+\dd y^2)/y^2$ 的曲率为 $-1$。最后一个模型是\emph{双曲平面（hyperbolic plane）}的一种表示，它首先是抽象黎曼度量，不应默认是整个上半平面在 $\R^3$ 中某个全局图像的诱导度量。

另一个例子是商空间 $\R^2/\mathbb Z^2$，把相差整数向量的点视为同一点。平移保持欧氏度量，因此可在商空间上定义平坦度量，得到紧致平坦环面。它与例~\ref{exam:dg-torus} 中的标准旋转环面具有相同拓扑，但后者曲率正负变化，两者不等距。拓扑类型、光滑结构、黎曼度量和欧氏空间中的放置方式是不同层次的数据。

后续学习可以沿两条相互联系的路线展开。一条路线研究流形、微分形式、积分与拓扑，理解为何曲率积分会产生拓扑不变量；另一条路线研究联络、测地线变分、Jacobi 场、曲率比较与整体黎曼几何。对曲面基本形式和测地线的熟练掌握，为这些抽象概念提供了可反复检验的具体模型。

\section{习题}
\begin{exercise}\label{ex:dg-manifolds-gradient}
在极坐标平面 $r>0$ 上求 $f(r,\theta)$ 的梯度，特别计算 $f=r^2$ 的梯度与 Laplace--Beltrami 算子。
\end{exercise}
\begin{exercise}\label{ex:dg-manifolds-sphere-laplace}
在半径 $R$ 球面上用经纬坐标计算 $f(\theta,\varphi)=\cos\theta$ 的 $\Delta f$，并解释极点处公式如何理解。
\end{exercise}
\begin{exercise}\label{ex:dg-manifolds-flat-torus}
说明欧氏度量为何能下降到 $\R^2/\mathbb Z^2$，并证明其测地线为直线投影；判断速度 $(a,b)\ne(0,0)$ 时测地线闭合的充要条件。
\end{exercise}
\begin{exercise}\label{ex:dg-manifolds-scale}
二维黎曼度量乘以正常数 $c^2$ 后，求长度、面积、高斯曲率与 $\int K\dd A$ 的变化规律。
\end{exercise}

